% ============================================================
% 03_abstract.tex — Résumés
% ============================================================

% ── RÉSUMÉ EN FRANÇAIS ───────────────────────────────────────
\newpage
\thispagestyle{empty}

\begin{center}
{\Large\bfseries\color{PrimaryBlue} Résumé}
\end{center}
\vspace{0.3cm}
\rule{\linewidth}{0.4pt}
\vspace{0.3cm}

{\normalsize
Face à la croissance rapide des agglomérations urbaines
marocaines et aux défis croissants de la congestion routière,
ce mémoire présente la conception et le déploiement d'un
système intelligent de gestion adaptative des feux de
circulation pour les villes intelligentes, en intégrant
des contraintes strictes de protection des données
personnelles conformément à la Loi 09-08.

Le système proposé repose sur un pipeline complet déployé
sur Raspberry Pi 5 — un dispositif embarqué à ressources
limitées — et s'articule autour de six modules interconnectés :
acquisition vidéo via OpenCV, détection multi-classes par
YOLO26n optimisé au format ONNX FP32, anonymisation temps réel
par floutage gaussien selon le principe Privacy-by-Design,
suivi des objets par tracker IoU sur zones polygonales,
décision adaptative des phases de feux via un Processus
de Décision Markovien (MDP) à politique déterministe,
et visualisation par un tableau de bord Flask en temps réel.

Un dataset hybride de \textbf{6 761 images} a été constitué
en trois étapes méthodologiques rigoureuses — split sans
fuite de données, renforcement de la classe \textit{person}
via CityPersons, et génération de crops pour les petits
piétons — à partir de sources internationales dont le
\textbf{Bellevue Traffic Video Dataset} et de données
locales collectées à \textbf{Lqliaa, Agadir}.
Un benchmark de six architectures YOLO (nano et small)
sur GPU et sur Raspberry Pi 5 confirme que \textbf{YOLO26n
FP32} atteint \textbf{7,6 FPS réels} sur Pi 5 — meilleur
compromis vitesse/taille parmi les douze configurations
testées. Les variantes small sont empiriquement confirmées
incompatibles avec le temps réel (moins de 3 FPS).

La politique MDP validée sur simulation SUMO démontre
une réduction du temps d'attente moyen de
\textbf{$-51\%$ à $-93\%$} selon le scénario,
par rapport aux feux à cycles fixes.
Les quatre hypothèses de recherche sont confirmées,
démontrant la viabilité d'un pipeline complet
d'intelligence artificielle embarquée pour la gestion
intelligente du trafic urbain, ouvrant la voie à des
déploiements réels dans les smart cities marocaines.
}

\vspace{0.4cm}

\begin{tcolorbox}[
  colback  = LightGray,
  colframe = DarkGray,
  arc      = 3pt,
  boxrule  = 0.5pt
]
\textbf{Mots-clés :} Edge AI, YOLO26n, Raspberry Pi 5,
Gestion du trafic, Smart City, Privacy-by-Design,
Anonymisation, ONNX FP32, Processus de Décision Markovien,
Loi 09-08, Bellevue Traffic Dataset.
\end{tcolorbox}

% ── ABSTRACT EN ANGLAIS ──────────────────────────────────────
\newpage
\thispagestyle{empty}

\begin{center}
{\Large\bfseries\color{PrimaryBlue} Abstract}
\end{center}
\vspace{0.3cm}
\rule{\linewidth}{0.4pt}
\vspace{0.3cm}

{\normalsize
In response to the rapid growth of Moroccan urban areas and
the increasing challenges of road congestion, this thesis
presents the design and deployment of an intelligent adaptive
traffic light management system for smart cities, integrating
strict personal data protection constraints in compliance
with Law 09-08.

The proposed system relies on a complete pipeline deployed
on a Raspberry Pi 5 — a resource-constrained embedded device
— structured around six interconnected modules: video
acquisition via OpenCV, multi-class detection using YOLO26n
in ONNX FP32 format, real-time anonymization through
Gaussian blurring following the Privacy-by-Design principle,
IoU-based object tracking over calibrated polygonal zones,
adaptive traffic phase decision-making via a deterministic
Markov Decision Process (MDP), and real-time visualization
through a Flask dashboard.

A hybrid dataset of \textbf{6,761 images} was built in
three methodological steps — leak-free group-based split,
person class reinforcement via CityPersons, and tiny
person crop generation — combining the
\textbf{Bellevue Traffic Video Dataset} with local data
collected in \textbf{Lqliaa, Agadir}.
A comprehensive benchmark of six YOLO architectures
(nano and small variants) on both GPU and Raspberry Pi 5
empirically confirms that \textbf{YOLO26n FP32} achieves
\textbf{7.6 real FPS} on Pi 5 — the best speed/size
trade-off among the twelve tested configurations.
Small variants are empirically confirmed incompatible
with real-time processing (below 3 FPS on Pi 5).

The MDP policy validated on SUMO simulation demonstrates
a reduction in average waiting time of
\textbf{$-51\%$ to $-93\%$} depending on the scenario,
compared to fixed-cycle traffic lights.
All four research hypotheses are confirmed, demonstrating
the viability of a complete embedded AI pipeline for
intelligent urban traffic management, paving the way
for real deployments in Moroccan smart cities.
}

\vspace{0.4cm}

\begin{tcolorbox}[
  colback  = LightGray,
  colframe = DarkGray,
  arc      = 3pt,
  boxrule  = 0.5pt
]
\textbf{Keywords:} Edge AI, YOLO26n, Raspberry Pi 5,
Traffic Management, Smart City, Privacy-by-Design,
Anonymization, ONNX FP32, Markov Decision Process,
Law 09-08, Bellevue Traffic Dataset.
\end{tcolorbox}

% ── RÉSUMÉ EN ARABE ──────────────────────────────────────────
\ifPDFTeX
\newpage
\thispagestyle{empty}
\begin{center}
\includegraphics[width=\textwidth,height=0.94\textheight,keepaspectratio]
  {generated/abstract_ar}
\end{center}
\else
\newpage
\thispagestyle{empty}

\begin{center}
{\Large\bfseries\color{PrimaryBlue}
\textarabic{ملخص}}
\end{center}
\vspace{0.3cm}
\rule{\linewidth}{0.4pt}
\vspace{0.3cm}

\begin{Arabic}

في ظل النمو المتسارع للتجمعات الحضرية المغربية
والتحديات المتزايدة لازدحام حركة المرور،
يقدم هذا البحث تصميم ونشر نظام ذكي
لإدارة إشارات المرور بشكل تكيفي في المدن الذكية،
مع دمج متطلبات صارمة لحماية البيانات الشخصية
وفقاً للقانون \LR{09-08}.

يعتمد النظام المقترح على خط معالجة متكامل
منشور على \LR{Raspberry Pi 5}،
ويتمحور حول ستة وحدات مترابطة:
استحواذ الفيديو عبر \LR{OpenCV}،
والكشف متعدد الفئات باستخدام \LR{YOLO26n}
بتنسيق \LR{ONNX FP32}،
وإخفاء الهوية في الوقت الفعلي
عبر التمويه الغاوسي وفق مبدأ
\LR{Privacy-by-Design}،
وتتبع الكائنات بواسطة خوارزمية \LR{IoU}
على مناطق متعددة الأضلاع،
واتخاذ قرار تكيفي لمراحل إشارات المرور
عبر عملية قرار ماركوف \LR{(MDP)}،
والتصور عبر لوحة تحكم \LR{Flask}.

تم بناء مجموعة بيانات هجينة من
\textbf{\LR{6\,761} صورة}
في ثلاث مراحل منهجية صارمة،
من مصادر دولية أبرزها
\LR{Bellevue Traffic Video Dataset}
ومن بيانات محلية جُمعت في
\textbf{لقليعة، أكادير}.
أثبت بنشمارك شامل لست معماريات \LR{YOLO}
على \LR{Raspberry Pi 5}
أن \LR{YOLO26n FP32} يحقق
\textbf{\LR{7.6} إطار في الثانية}،
أفضل توازن بين السرعة والحجم
من بين اثنتي عشرة تكويناً مختبراً.

أثبتت سياسة \LR{MDP} المُتحقق منها
عبر محاكاة \LR{SUMO} تخفيضاً
في متوسط وقت الانتظار من
\textbf{\LR{51\%}-} إلى \textbf{\LR{93\%}-}
مقارنةً بإشارات المرور ذات الدورات الثابتة.
تم التحقق من الفرضيات الأربع للبحث،
مما يُثبت جدوى خط معالجة متكامل
للذكاء الاصطناعي المدمج لإدارة
حركة المرور الحضرية بذكاء،
فاتحاً الطريق أمام عمليات النشر
الفعلية في المدن الذكية المغربية.

\end{Arabic}

\vspace{0.4cm}

\begin{tcolorbox}[
  colback  = LightGray,
  colframe = DarkGray,
  arc      = 3pt,
  boxrule  = 0.5pt
]
\begin{Arabic}
\textbf{الكلمات المفتاحية:}
الذكاء الاصطناعي المدمج،
\LR{YOLO26n}،
\LR{Raspberry Pi 5}،
إدارة حركة المرور،
المدينة الذكية،
الخصوصية بالتصميم،
إخفاء الهوية،
\LR{ONNX FP32}،
عملية قرار ماركوف،
القانون \LR{09-08}،
\LR{Bellevue Traffic Dataset}.
\end{Arabic}
\end{tcolorbox}
\fi